{{دیوان
| عنوان = تہذیبِ حاضر
| ماخذ = ویکی ماخذ
| ماخذ_لنک = https://ur.wikisource.org/wiki/%D8%AA%DB%81%D8%B0%DB%8C%D8%A8%D9%90_%D8%AD%D8%A7%D8%B6%D8%B1
| مقابلہ = اقبال سائبر لائبریری (اقبال اکادمی پاکستان)
| مقابلہ_لنک = https://www.iqbalcyberlibrary.net
}}
=== تضمین بر شعرِ فیضیؔ ===

<poem>
حرارت ہے بلا کی بادۂ تہذیبِ حاضر میں
بھڑک اُٹھّا بھبُوکا بن کے مُسلم کا تنِ خاکی

کِیا ذرّے کو جُگنو دے کے تابِ مُستعار اس نے
کوئی دیکھے تو شوخی آفتابِ جلوہ فرما کی

نئے انداز پائے نوجوانوں کی طبیعت نے
یہ رعنائی، یہ بیداری، یہ آزادی، یہ بے باکی

تغیّر آگیا ایسا تدبّر میں، تخیّل میں
ہنسی سمجھی گئی گُلشن میں غُنچوں کی جگر چاکی

کِیا گُم تازہ پروازوں نے اپنا آشیاں لیکن
مناظر دِلکُشا دِکھلا گئی ساحر کی چالاکی

حیاتِ تازہ اپنے ساتھ لائی لذّتیں کیا کیا
رقابت، خودفروشی، ناشکیبائی، ہوسناکی

فروغِ شمعِ نَو سے بزمِ مسلم جگمگا اُٹھّی
مگر کہتی ہے پروانوں سے میری کُہنہ اِدراکی

“تو اے پروانہ! ایں گرمی ز شمع محفلے داری
چو من در آتشِ خود سوز اگر سوزِ دلے داری”
</poem>
